\chapter[\emj{↩️}~Backtracking (Vuelta Atrás)]{\emj{↩️}~Backtracking (Vuelta Atrás)}

\begin{dsaquees}
Resolver un laberinto 🌀 con un lápiz. 

Sigues un camino con la punta del lápiz hasta que llegas a un 
callejón sin salida. En lugar de rendirte, borras lo que hiciste 
(vuelta atrás) hasta la última bifurcación y pruebas el otro camino 
que dejaste pendiente.

Backtracking es probar TODAS las posibilidades una por una, pero 
con la inteligencia de "deshacer" tus pasos cuando te das cuenta 
de que una ruta no lleva a la solución. Es una búsqueda exhaustiva 
pero controlada.
\end{dsaquees}
\begin{dsaotraforma}
Backtracking es una técnica algorítmica para resolver problemas de 
forma recursiva, intentando construir una solución incrementalmente. 

Se basa en la exploración de un "Árbol de Decisiones". Si en algún 
punto la solución parcial no cumple con las restricciones del problema 
(se vuelve inválida), el algoritmo descarta esa rama ("pruning") y 
regresa al estado anterior para intentar otra opción.

Es la base para resolver problemas de combinatoria, permutaciones y 
juegos como el Sudoku o el Ajedrez.
\end{dsaotraforma}
\begin{dsadiagrama}
[root]
              /      \
        Elegir 1?   NO elegir 1?
         /    \        /      \
     Elegir 2? NO?   Elegir 2? NO?
       |        |       |        |
     [1,2]     [1]     [2]      []

¡Cada nodo es una decisión! Backtracking recorre este árbol usando DFS.
\end{dsadiagrama}
\begin{dsacomofunciona}
Casi todos los problemas de Backtracking siguen este patrón:

1. Caso Base: ¿Ya terminamos de construir la solución?
   - SÍ: Guárdala o imprímela. Regresa.
2. Opciones: Para cada opción disponible en este nivel:
   - ¿Es válida la opción? (Restricciones).
   - SÍ: 
     a. Hacer Elección: Modifica el estado actual.
     b. Recursión: Llama a la función para el siguiente nivel.
     c. Deshacer Elección (BACKTRACK): Regresa el estado como estaba.
\end{dsacomofunciona}
\section*{PSEUDOCÓDIGO}
\hrulefill
\begin{verbatim}
función backtrack(estado, opciones):
    SI estado es meta:
        imprimir(estado)
        return

    PARA cada o EN opciones:
        SI esValida(o, estado):
            agregar o a estado // "Hacer"
            backtrack(estado, opciones\_restantes)
            quitar o de estado // "Deshacer" (La clave ✨)
\end{verbatim}
\section*{CÓDIGO C++}
\hrulefill
\begin{lstlisting}
#include <iostream>
#include <vector>
#include <string>
using namespace std;

// Problema: Generar todas las permutaciones de un string -> O(n!)
void permute(string& s, int l, int r) {
    // Caso Base
    if (l == r) {
        cout << s << endl;
        return;
    }

    for (int i = l; i <= r; i++) {
        // 1. Hacer Elección (Swap)
        swap(s[l], s[i]);

        // 2. Recursión (Ir al siguiente nivel)
        permute(s, l + 1, r);

        // 3. Deshacer Elección (Backtrack)
        // Regresamos el string a su estado original para la sig. iteración
        swap(s[l], s[i]);
    }
}

int main() {
    string str = "ABC";
    cout << "Permutaciones de ABC:" << endl;
    permute(str, 0, str.length() - 1);
    return 0;
}
\end{lstlisting}
\begin{dsacomplejidad}
Problema        │ Tiempo Típico │ Razón
  ────────────────┼───────────────┼──────────────────────────────
  Permutaciones   │ O(n!)         │ n * (n-1) * (n-2)...
  Subconjuntos    │ O(2\textasciicircum{}n)        │ Cada elemento se elige o no
  N-Queens        │ O(n!)         │ Menos gracias al pruning
  Sudoku          │ O(9\textasciicircum{}(n*n))    │ Fuerza bruta pura (peor caso)

* El espacio suele ser O(n) por la profundidad del stack de recursión.
\end{dsacomplejidad}
\begin{dsaentrevistas}

\end{dsaentrevistas}
\begin{dsaresumen}
• DFS + Deshacer cambios.
• Explora todas las soluciones posibles.
• Úsalo para permutaciones, combinaciones y laberintos.
• Complejidad alta (Exponencial/Factorial).
• Pruning: Cortar ramas que no sirven para ganar velocidad.
• Esencial dominar el template: Hacer -> Recursión -> Deshacer.
════════════════════════════════════════════════════════════════
\end{dsaresumen}

